% Источник: Конспект
\ticket{05}{Лемма об условиях, эквивалентных непрерывности линейного оператора}
\label{ticket:05}

\begin{definition}
Линейный оператор $A:X\to Y$ между нормированными пространствами называется \textbf{ограниченным}, если $\|A\|<+\infty$.
Эквивалентно, существует $C<\infty$ такое, что $\|Ax\|_Y\le C\|x\|_X$ для всех $x\in X$.
Он \textbf{равномерно непрерывен}, если для каждого $\varepsilon>0$ найдётся $\delta>0$ такое, что
\[
  \|x-y\|_X<\delta\quad\Longrightarrow\quad
  \|Ax-Ay\|_Y<\varepsilon\qquad(x,y\in X).
\]
\end{definition}

\begin{lemma}
Пусть $X$, $Y$ --- нормированные пространства, $A:X\to Y$ --- линейный оператор.
Тогда эквивалентны:
\begin{enumerate}
  \item $A$ ограничен;
  \item $A$ непрерывен в нуле;
  \item $A$ непрерывен на $X$;
  \item $A$ равномерно непрерывен на $X$.
\end{enumerate}
\end{lemma}

\begin{proof}
Переходы \(4 \Rightarrow 3 \Rightarrow 2\) очевидны.

\medskip
\noindent\textbf{\(2 \Rightarrow 1\).}
Для \(\varepsilon=1\) существует \(\delta>0\), такое что
\[
    \forall x\in B(0,\delta)
    \qquad |Ax|<1.
\]

Для любого вектора \(u\), такого что \(|u|=1\), имеем
\[
    \left|\frac{\delta}{2}u\right|
    =\frac{\delta}{2}<\delta.
\]
Поэтому
\[
    \left|A\left(\frac{\delta}{2}u\right)\right|<1.
\]
Из линейности \(A\) следует
\[
    |Au|
    =\frac{2}{\delta}
      \left|A\left(\frac{\delta}{2}u\right)\right|
    <\frac{2}{\delta}.
\]
Следовательно, \(A\) ограничен, причём
\[
    \|A\|\le\frac{2}{\delta}.
\]

\medskip
\noindent\textbf{\(1 \Rightarrow 4\).}
Если \(A=0\), утверждение очевидно.
Пусть \(\|A\|>0\).

Для любого \(\varepsilon>0\) положим
\[
    \delta=\frac{\varepsilon}{\|A\|}.
\]
Тогда для любых \(x,\bar{x}\), удовлетворяющих условию
\[
    |x-\bar{x}|<\delta,
\]
получаем
\[
    |Ax-A\bar{x}|
    =|A(x-\bar{x})|
    \le\|A\|\cdot|x-\bar{x}|
    <\|A\|\delta
    =\varepsilon.
\]
Таким образом, \(A\) равномерно непрерывен.
\end{proof}
